\section{Robustez e inferencia}\label{sec:robustness}

En esta sección someto los resultados principales a una batería de pruebas de
especificación, de inferencia y de diseño, y luego a una replicación fuera de muestra. El
mensaje general es que los resultados de composición, la caída de la formalidad y el
aumento del trabajo independiente, son los hallazgos más robustos, y los corrobora un
diseño independiente de exposición continua con inferencia por aleatorización; la
evidencia salarial muestra de forma consistente que no hubo ganancias de remuneración,
aunque con una salvedad sobre las tendencias previas; y el resultado sobre el nivel de
empleo no supera sus diagnósticos y no debe leerse como causal.

\subsection{Robustez a la especificación}

El Cuadro~\ref{tab:robspec} reporta el coeficiente $\mathrm{Low}\times\mathrm{Post}$ para
los cuatro resultados principales bajo diez especificaciones alternativas: sin controles;
sin el trimestre pandémico de inicios de 2021; con tendencias lineales específicas por
departamento; restringiendo la muestra a trabajadores en edad central; con un corte
educativo alternativo que elimina las categorías situadas en la frontera entre ambos
grupos; con efectos fijos de ocupación y de rama de actividad; agrupando los errores por
departamento y nivel de calificación; reincorporando el trimestre de transición 2022Q2,
parcialmente tratado; excluyendo la agricultura (cubierta por el régimen laboral agrario
de la Ley 31110 y no por el salario mínimo general); y reincorporando a los trabajadores
del sector público en las muestras condicionadas a tener empleo. Las estimaciones de
formalidad y de trabajo independiente son estables en signo y magnitud en todos los
casos y se mueven dentro de un rango estrecho. La estimación salarial es negativa o nula
en todas las especificaciones y nunca es significativamente positiva. La estimación de
empleo, en cambio, se atenúa notablemente cuando se excluye el trimestre pandémico o se
incluyen tendencias específicas por departamento, lo que refuerza mi renuencia a
interpretarla de manera causal.

\begin{table}[!tbp]\centering
\caption{Robustez de las estimaciones DiD a la especificación}\label{tab:robspec}
\begin{threeparttable}
\footnotesize
\setlength{\tabcolsep}{4pt}
\begin{tabular}{lcccc}
\toprule
Especificación & Log salario/hora & Empleo & Empleo formal & Trab. independiente \\
\midrule
Base & -0.021 & -0.025$^{***}$ & -0.046$^{***}$ & 0.036$^{***}$ \\
 & (0.012) & (0.006) & (0.007) & (0.006) \\
 & \footnotesize[69\,534] & \footnotesize[278\,064] & \footnotesize[181\,317] & \footnotesize[181\,317] \\
Sin controles & -0.008 & -0.021$^{***}$ & -0.042$^{***}$ & 0.042$^{***}$ \\
 & (0.012) & (0.006) & (0.007) & (0.006) \\
 & \footnotesize[69\,534] & \footnotesize[278\,064] & \footnotesize[181\,317] & \footnotesize[181\,317] \\
Sin 2021Q1 (pandemia) & -0.023$^{*}$ & -0.014$^{***}$ & -0.042$^{***}$ & 0.031$^{***}$ \\
 & (0.012) & (0.005) & (0.008) & (0.005) \\
 & \footnotesize[65\,279] & \footnotesize[259\,135] & \footnotesize[168\,908] & \footnotesize[168\,908] \\
Tendencias lineales por depto. & -0.022$^{*}$ & -0.017$^{***}$ & -0.040$^{***}$ & 0.034$^{***}$ \\
 & (0.012) & (0.006) & (0.007) & (0.006) \\
 & \footnotesize[69\,534] & \footnotesize[278\,064] & \footnotesize[181\,317] & \footnotesize[181\,317] \\
Edad central (25--55) & -0.030$^{*}$ & -0.023$^{***}$ & -0.042$^{***}$ & 0.029$^{***}$ \\
 & (0.015) & (0.004) & (0.008) & (0.007) \\
 & \footnotesize[46\,021] & \footnotesize[161\,115] & \footnotesize[119\,540] & \footnotesize[119\,540] \\
Corte educativo alternativo & -0.015 & -0.039$^{***}$ & -0.065$^{***}$ & 0.050$^{***}$ \\
 & (0.012) & (0.007) & (0.006) & (0.007) \\
 & \footnotesize[40\,526] & \footnotesize[186\,263] & \footnotesize[117\,174] & \footnotesize[117\,174] \\
EF de ocupación y actividad & -0.007 & -0.016$^{***}$ & -0.036$^{***}$ & 0.027$^{***}$ \\
 & (0.011) & (0.003) & (0.006) & (0.005) \\
 & \footnotesize[69\,534] & \footnotesize[213\,810] & \footnotesize[181\,317] & \footnotesize[181\,317] \\
Agrupación dpto.$\times$ed. & -0.021 & -0.025 & -0.046$^{***}$ & 0.036$^{***}$ \\
 & (0.013) & (0.020) & (0.007) & (0.007) \\
 & \footnotesize[69\,534] & \footnotesize[278\,064] & \footnotesize[181\,317] & \footnotesize[181\,317] \\
Tendencia lineal, baja educación & -0.021 & -0.005 & -0.004 & 0.006 \\
 & (0.019) & (0.008) & (0.010) & (0.009) \\
 & \footnotesize[69\,534] & \footnotesize[278\,064] & \footnotesize[181\,317] & \footnotesize[181\,317] \\
Con trim. de transición (2022Q2) & -0.018 & -0.025$^{***}$ & -0.044$^{***}$ & 0.035$^{***}$ \\
 & (0.012) & (0.005) & (0.007) & (0.006) \\
 & \footnotesize[74\,310] & \footnotesize[296\,954] & \footnotesize[193\,762] & \footnotesize[193\,762] \\
Sin agricultura (régimen agrario) & -0.022$^{*}$ & -- & -0.038$^{***}$ & 0.028$^{***}$ \\
 & (0.012) &  & (0.007) & (0.005) \\
 & \footnotesize[53\,755] & \footnotesize[--] & \footnotesize[114\,212] & \footnotesize[114\,212] \\
Con sector público & -0.002 & -- & -0.037$^{***}$ & 0.030$^{***}$ \\
 & (0.010) &  & (0.008) & (0.006) \\
 & \footnotesize[87\,033] & \footnotesize[--] & \footnotesize[200\,383] & \footnotesize[200\,383] \\
\bottomrule
\end{tabular}
\begin{tablenotes}\footnotesize
\item Notas: Cada celda es el coeficiente DiD de $\mathrm{Low}\times\mathrm{Post}$ (debajo, el error estándar agrupado por departamento; entre corchetes, el tamaño de muestra) para el resultado del encabezado de la columna, bajo la especificación de la fila; todas las estimaciones están ponderadas con los factores de expansión de la ENAHO, y Log salario por hora denota el log del salario real por hora. La especificación de base es la ecuación~\eqref{eq:twfe}, que excluye el trimestre 2022Q2, parcialmente tratado. El corte educativo alternativo define baja educación como a lo sumo secundaria incompleta y alta educación como al menos superior no universitaria completa, y excluye las categorías límite. La exclusión de la agricultura y la inclusión del sector público se aplican solo a los resultados condicionados al empleo (la actividad y el sector no están definidos para quienes no trabajan). $^{*}p<0.1$, $^{**}p<0.05$, $^{***}p<0.01$.
\end{tablenotes}
\end{threeparttable}\end{table}


\subsection{Inferencia}

Como el contraste de tratamiento varía solo entre dos grupos educativos, el
agrupamiento convencional por departamento no proporciona por sí solo una inferencia
basada en el diseño para el coeficiente $\mathrm{Low}\times\mathrm{Post}$; por ello
reporto un conjunto de procedimientos cada vez más exigentes. El
Cuadro~\ref{tab:inference} reúne la evidencia. Las columnas (1) a (3) reportan los
valores $p$ del efecto base con errores agrupados por departamento (25 conglomerados),
por departamento y nivel de calificación (50 conglomerados), y con la corrección CR2
para muestras pequeñas de \citet{pustejovsky_tipton_2018}, aplicada a celdas colapsadas
en el espíritu de \citet{bertrand_2004}, aquí al nivel
departamento$\times$grupo$\times$periodo, que es la agregación más fina en la que la
corrección CR2 resulta computacionalmente factible.
Los efectos sobre la formalidad y el trabajo independiente siguen siendo altamente
significativos bajo ambos esquemas de agrupamiento; con el procedimiento CR2 sobre
celdas colapsadas, deliberadamente conservador y de baja potencia con veinticinco
departamentos, el trabajo independiente sigue siendo marginalmente significativo y la
formalidad no. El efecto sobre el empleo pierde significancia con cualquier
procedimiento más exigente que el simple agrupamiento por departamento. Como una prueba
de permutación del contraste entre dos grupos no está definida, la inferencia por
aleatorización basada en el diseño se aplica, en su lugar, al diseño de exposición
continua, donde la asignación que se permuta (la incidencia regional) toma veinticinco
valores: la columna (4) reporta estos valores $p$, que corresponden a las estimaciones
de la columna (5) y no al contraste educativo.\footnote{Un \emph{wild cluster
bootstrap} en el espíritu de \citet{cameron_2008_bootstrap} no está disponible para
estimaciones ponderadas con factores de expansión, de modo que los procedimientos CR2 y
de aleatorización cumplen aquí ese papel.}

\begin{table}[!tbp]\centering
\caption{Robustez de la inferencia, robustez del diseño y validación fuera de muestra}\label{tab:inference}
\begin{threeparttable}
\footnotesize
\setlength{\tabcolsep}{3.5pt}
\begin{tabular}{lcccccc}
\toprule
Resultado & (1) Dpto. & (2) Dpto.$\times$ed. & (3) CR2 & (4) RI (cont.) & (5) Continuo & (6) 2025 \\
\midrule
Log salario/hora & 0.106 & 0.126 & 0.854 & 0.029 & 0.017 & 0.002 \\
 & & & & & (0.012) & (0.011) \\
Empleo & 0.000 & 0.210 & 0.507 & 0.003 & -0.034$^{***}$ & -0.002 \\
 & & & & & (0.012) & (0.005) \\
Empleo formal & 0.000 & 0.000 & 0.244 & 0.000 & -0.018$^{***}$ & -0.001 \\
 & & & & & (0.003) & (0.008) \\
Trab. independiente & 0.000 & 0.000 & 0.093 & 0.024 & 0.010$^{**}$ & -0.014$^{**}$ \\
 & & & & & (0.004) & (0.006) \\
\bottomrule
\end{tabular}
\begin{tablenotes}\footnotesize
\item Notas: Las columnas (1)--(3) reportan valores $p$ del efecto de base $\mathrm{Low}\times\mathrm{Post}$ bajo, respectivamente, agrupación por departamento (25 grupos), agrupación por departamento$\times$educación (50 grupos) y la corrección CR2 para muestras pequeñas \citep{pustejovsky_tipton_2018} aplicada a celdas colapsadas a nivel de departamento. La columna (4) reporta el valor $p$ de inferencia por aleatorización (RI) del diseño de exposición CONTINUA de la columna (5), no del contraste de dos grupos $\mathrm{Low}\times\mathrm{Post}$, obtenido permutando los 25 índices regionales de Kaitz entre departamentos (2000 extracciones). La columna (5) reporta la estimación de exposición continua, el coeficiente de $\mathrm{Post}\times$(índice de Kaitz departamental estandarizado), con errores estándar agrupados por departamento entre paréntesis. La columna (6) reporta la estimación DiD de $\mathrm{Low}\times\mathrm{Post}$ para la reforma de enero de 2025 (ventana 2024--2025), con errores estándar agrupados por departamento entre paréntesis. Log salario por hora denota el log del salario real por hora; todas las estimaciones están ponderadas con los factores de expansión de la ENAHO. $^{*}p<0.1$, $^{**}p<0.05$, $^{***}p<0.01$.
\end{tablenotes}
\end{threeparttable}\end{table}


\subsection{Robustez del diseño: intensidad de la incidencia y triples diferencias}

Mi identificación descansa en la exposición diferencial por nivel educativo, pero el
salario mínimo también tiene una incidencia (\emph{bite}) distinta entre regiones, y esto
ofrece una fuente independiente de variación que no depende en absoluto del contraste
educativo. La columna (5) del Cuadro~\ref{tab:inference} reporta una especificación de
exposición continua en el espíritu de \citet{card_1992}, que relaciona los resultados
con el índice de Kaitz departamental previo a la reforma, estandarizado e interactuado
con el indicador del periodo posterior. Los resultados corroboran el diseño basado en la
calificación en los márgenes de composición: en los departamentos donde el salario
mínimo incide con más fuerza, tras la reforma el empleo formal cae 1.8 puntos
porcentuales más por cada desviación estándar de exposición ($p<0.001$; $p$ por
aleatorización $<0.001$) y el trabajo independiente aumenta 1.0 punto más ($p=0.02$;
$p$ por aleatorización $=0.02$). En este diseño el empleo también disminuye con la
exposición, aunque descarto ese margen por las razones ya expuestas. Los salarios no
muestran una respuesta significativa. Una especificación de triple diferencia que
interactúa el contraste por calificación con un indicador de departamentos de alta
incidencia no es estadísticamente significativa para ningún resultado (Cuadro del
Apéndice~\ref{tab:ddd}); lo interpreto como falta de potencia de una partición regional
binaria y gruesa, más que como evidencia contra el mecanismo, dado que la versión
continua es informativa y que las estimaciones por subgrupos de la
Sección~\ref{sec:heterogeneity} se mueven en la dirección esperada.

Otras tres pruebas de diseño matizan esta corroboración (Cuadros del
Apéndice~\ref{tab:expocells} y~\ref{tab:sectorocc}). En primer lugar, el diseño regional
supera su propio diagnóstico de adelantos justo donde importa: las pruebas conjuntas de
las interacciones trimestre$\times$exposición previas a la reforma no rechazan la nula
para la formalidad ($p=0.84$), el trabajo independiente ($p=0.20$) ni los salarios
($p=0.24$), y solo la rechazan para el empleo ($p=0.002$), lo que es coherente con todo
lo demás que sabemos de ese margen. En segundo lugar, en aras de la plena
transparencia, estimo también diseños más finos de fracción afectada a nivel de celda
\citep{card_1992, clemens_wither_2019}, con la exposición medida dentro de celdas
departamento$\times$calificación$\times$edad; ambas variantes fallan sus diagnósticos de
adelantos: la exposición de grano fino se confunde con las heterogéneas recuperaciones
pospandemia de las celdas, y la variante entre pisos incluso invierte los signos, porque
las celdas con masa entre ambos pisos son precisamente las celdas semiformales que se
recuperaron más rápido. Por ello, trato al departamento como el nivel creíble para la
variación de la exposición y reporto los diseños por celda en el apéndice en lugar de
omitirlos. En tercer lugar, la reasignación es amplia entre sectores: la caída de la
formalidad aparece por igual en los sectores cubiertos de alta incidencia, en los demás
sectores privados y en la agricultura. La estimación agrícola no es una falla de
falsificación, porque la remuneración diaria del régimen agrario está indexada a la RMV
y, por tanto, subió con la reforma; en el Perú ningún sector grande queda con su piso
sin cambios, y por eso los contrastes de cobertura no pueden ser nítidos en este caso.
Una relación dosis-respuesta a nivel de ocupación (con la exposición medida como la
proporción de asalariados de la ocupación que ganaban menos que el nuevo piso antes de
la reforma) muestra mayores caídas en la remuneración por debajo del piso y, de manera
sugerente, ganancias salariales relativas en las ocupaciones más expuestas, aunque la
ocupación se mide después del tratamiento y leo ese panel como descriptivo.

Como prueba final de diseño, estimo diferencias en diferencias sintéticas
\citep{arkhangelsky_2021_sdid} sobre celdas departamento-trimestre, tratando a los
departamentos de alta incidencia como tratados desde 2022Q3 y dejando que el estimador
repondere a los departamentos de baja incidencia para reconstruir su trayectoria previa
a la reforma (Cuadro del Apéndice~\ref{tab:sdidpower}). Las estimaciones van en la
dirección de la reasignación (la formalidad cae y el trabajo independiente sube, ya sea
que el resultado sea la media de baja educación o la brecha entre baja y alta
calificación), con magnitudes de aproximadamente la mitad de las estimaciones TWFE y una
precisión limitada por las veinticinco unidades agregadas; la caída de la formalidad en
las medias de baja educación es significativa al cinco por ciento con el
\emph{jackknife}.

\subsection{Selección y potencia}

Como la reforma sacó a trabajadores de la muestra de asalariados, las diferencias en
diferencias salariales se condicionan a una población endógena. El Cuadro del
Apéndice~\ref{tab:leemde} acota el sesgo resultante con el método de recorte de
\citet{lee_2009}: la reforma redujo la inclusión de los trabajadores poco calificados en
la muestra salarial en 2.1 puntos porcentuales (7.6 por ciento de la tasa base), y
recortar en esa fracción la distribución salarial low$\times$post desde cualquiera de
sus colas acota las diferencias en diferencias salariales 2$\times$2 no condicionadas
entre $-0.10$ y $+0.10$. La selección por sí sola podría, entonces, ocultar efectos
salariales de hasta diez puntos logarítmicos en cualquier dirección, lo cual es una
declaración franca de lo que la evidencia salarial puede y no puede descartar; la
evidencia de cumplimiento y de acumulación (\emph{bunching}) de la
Sección~\ref{sec:results}, que solo condiciona en la condición de formalidad, no está
sujeta a esta preocupación, y es en ella donde apoyo la conclusión sobre la
fiscalización. El mismo cuadro del apéndice reporta los efectos mínimos detectables: con
una potencia del ochenta por ciento, el diseño puede detectar efectos salariales de 3.4
puntos logarítmicos y efectos de 2.1 y 1.7 puntos porcentuales sobre la formalidad y el
trabajo independiente, ambos muy por debajo de los estimados de 4.6 y 3.6, de modo que
los hallazgos de composición no son un artefacto de una baja potencia estadística.

La misma lógica de selección se aplica a los propios resultados de composición, que se
definen entre los ocupados, mientras que el margen de empleo arrastra una tendencia
previa ligada a la recuperación. Por ello, el Cuadro del Apéndice~\ref{tab:uncond}
reporta versiones no condicionadas, definidas como proporciones de la población en edad
de trabajar, que no condicionan en el empleo. En el contraste educativo, las
estimaciones no condicionadas son mayores (empleo formal $-3.3$ puntos, trabajo
independiente $+1.7$), pero heredan la tendencia previa diferencial de recuperación
pospandemia del margen de empleo y no superan su diagnóstico, de modo que las reporto
por transparencia, sin apoyarme en ellas. El resultado informativo está en el diseño
regional, que no usa el contraste educativo: el empleo formal no condicionado cae 1.6
puntos porcentuales por cada desviación estándar de exposición ($p$ por aleatorización
$<0.001$), de modo que la caída del empleo formal no es un artefacto de condicionar en
una población ocupada cambiante. El trabajo independiente no condicionado, en cambio, no
aumenta con la exposición regional; en ese diseño, el aumento de la \emph{proporción}
del trabajo independiente en el empleo opera a través de la reducción del denominador de
empleo, y la reasignación absoluta hacia el trabajo independiente es un rasgo propio del
contraste educativo. Retomo esta distinción en la Sección~\ref{sec:discussion}.

\subsection{Sensibilidad a violaciones de las tendencias paralelas}

Incluso unas tendencias previas planas se estiman con error, por lo que aplico el
análisis de sensibilidad de \citet{rambachan_roth_2023}. La Figura~\ref{fig:honest}
muestra el conjunto de confianza robusto para el efecto promedio posterior a la reforma
sobre el empleo y la formalidad, cuando permito que la violación posterior de las
tendencias paralelas sea tan grande como un múltiplo $\bar{M}$ de la mayor violación
previa a la reforma. Los efectos sobre la formalidad y el trabajo independiente se
mantienen alejados de cero bajo tendencias paralelas exactas, pero incluyen el cero en
cuanto se admiten violaciones de apenas $\bar{M}=0.1$, y el efecto sobre el empleo
incluye el cero incluso con $\bar{M}=0$. Este es el eslabón más débil del diseño
basado en la educación, y lo digo con claridad: los resultados de composición son
robustos a la inferencia convencional y a la basada en el diseño, pero no a violaciones
moderadas de las tendencias paralelas en el contraste educativo. Su credibilidad
descansa, por lo tanto, en las dos piezas de evidencia que no dependen de ese contraste:
el gradiente regional de dosis-respuesta y el diseño de exposición continua con
inferencia por aleatorización. El cálculo de potencia de \citet{roth_2022_pretest}
cuantifica por qué no me apoyo solo en haber superado las pruebas previas: una violación
lineal de las tendencias paralelas que mi prueba previa detectaría solo la mitad de las
veces sesgaría el efecto promedio sobre la formalidad en el periodo posterior en 0.042,
nueve décimas del efecto estimado (Cuadro del Apéndice~\ref{tab:sdidpower}, Panel B).
Superar las pruebas previas da tranquilidad, pero no es concluyente, y así las trato a
lo largo de todo el artículo.

La granularidad temporal tampoco determina las estimaciones. Los estudios de eventos
mensuales (Figura~\ref{fig:monthly}), que fechan los resultados salariales según el mes
de referencia del ingreso y los de empleo según el mes de la entrevista, muestran que la
caída de la formalidad aparece inmediatamente después de mayo de 2022 y persiste, con
una caída sugerente en las semanas en torno a la publicación del decreto, el 3 de abril;
y un estudio de eventos trimestral que agrupa la cola derecha en $k\ge 8$, en lugar de
estimar once rezagos libres, deja las estimaciones sin cambios.\footnote{Las series
completas de coeficientes mensuales y agrupados se incluyen en el paquete de
replicación.}

\begin{figure}[!t]
\centering
\includegraphics[width=0.98\textwidth]{fig11_monthly_event.pdf}
\caption{Estudios de eventos mensuales en torno a la reforma de mayo de 2022. Los
resultados salariales se fechan según el mes de referencia del ingreso y los de empleo
según el mes de la entrevista; extremos agrupados en $\pm$13 meses; mes de referencia:
abril de 2022. Bandas: intervalos de confianza al 95 por ciento con errores agrupados
por departamento.}
\label{fig:monthly}
\end{figure}

\begin{figure}[!t]
\centering
\includegraphics[width=0.8\textwidth]{fig7_honestdid.pdf}
\caption{Sensibilidad a violaciones de las tendencias paralelas
\citep{rambachan_roth_2023}. Conjuntos de confianza robustos para el efecto promedio
posterior a la reforma bajo la restricción de magnitudes relativas, que acota la
violación de la tendencia posterior a la reforma en $\bar{M}$ veces la mayor violación
previa a la reforma.}
\label{fig:honest}
\end{figure}

\subsection{Reformas placebo y exclusión de una unidad a la vez}

Asigno reformas ficticias a cada trimestre interior de la ventana previa a la reforma y
vuelvo a estimar las diferencias en diferencias usando solo datos previos a la reforma
(Cuadro del Apéndice~\ref{tab:placebo}). Los placebos salariales son nulos y los del
trabajo independiente son marginales en dos de los tres trimestres. Los placebos de
empleo son grandes y significativos en los tres trimestres, lo cual es un síntoma
directo de la tendencia previa impulsada por la pandemia en ese resultado y la base de
mi negativa a interpretar causalmente la estimación de empleo.

Los placebos de formalidad merecen algo más que una lectura de aprobado o desaprobado,
porque ellos y la prueba conjunta de tendencias previas responden la misma pregunta con
distinta potencia. Los tres son negativos ($-0.024$, $-0.029$, $-0.017$, el último
significativo al cinco por ciento): aunque la prueba conjunta no rechaza la nula
($p=0.15$), el patrón de las estimaciones puntuales indica una deriva relativa a la baja
de aproximadamente dos puntos porcentuales a lo largo de la ventana previa. Por eso
cuantifico cómo cambia la estimación de formalidad bajo correcciones de la deriva de
severidad creciente, en lugar de apoyarme solo en haber superado la prueba. Descontar
aritméticamente la deriva implícita en los placebos deja un efecto de aproximadamente
$-0.02$ a $-0.03$, que sigue siendo más de una décima parte de la formalidad base y que
está muy cerca de la estimación de la brecha por calificación con diferencias en
diferencias sintéticas ($-0.028$), inmune por construcción a una deriva a nivel de
grupo. Dos correcciones más agresivas acotan los extremos. Una especificación que
permite que los dos grupos educativos diverjan según una tendencia lineal libre a lo
largo de toda la ventana (Cuadro~\ref{tab:robspec}) lleva a cero las estimaciones de
composición del contraste educativo, pero esto es mecánico y no diagnóstico: un cambio
de nivel inmediato y persistente, observado durante diez trimestres posteriores a la
reforma, se aproxima bien mediante una tendencia, de modo que esta especificación
absorbe un verdadero efecto escalón junto con cualquier deriva y conviene leerla como
el peor escenario. De manera simétrica, la restricción de suavidad de
\citet{rambachan_roth_2023}, que extrapola la tendencia previa lineal ajustada al
periodo posterior, produce un intervalo no informativo para la formalidad
($[-0.045,\,0.041]$ con $M=0$): extrapolar cuatro puntos ruidosos del periodo previo a
lo largo de diez trimestres posteriores dispara la varianza. El resumen justo es que el
contraste educativo, por sí solo, no puede separar un cambio de nivel persistente de una
tendencia de grupo en una ventana tan larga. Precisamente por eso el peso causal de este
artículo descansa en el diseño de exposición regional, cuyos adelantos son planos, cuya
inferencia se basa en el diseño y cuyo gradiente de dosis-respuesta no puede ser
generado por una deriva específica de grupo; y en la convergencia de la estimación
educativa ajustada por la deriva con la estimación de diferencias en diferencias
sintéticas, de $-0.02$ a $-0.03$.

La reestimación excluyendo un departamento a la vez y excluyendo un trimestre a la vez
(Cuadro del Apéndice~\ref{tab:loo}) deja las estimaciones puntuales prácticamente sin
cambios para los cuatro resultados, lo que indica que ninguna región ni ningún
trimestre en particular determina los resultados.

\subsection{La ola de inmigración venezolana}

Un choque distinto de oferta laboral se superpone a mi contexto: el éxodo venezolano
trajo al Perú alrededor de 1.5 millones de migrantes, concentrados en 2018--2019 y en
Lima y la costa urbana, que compiten sobre todo por empleos urbanos informales y de
bajos salarios \citep{asencios_castellares_2020}. Tres rasgos de la evidencia indican que
este flujo no explica el resultado de reasignación (Cuadro del
Apéndice~\ref{tab:migration}). En primer lugar, añadir efectos fijos
departamento$\times$trimestre, que absorben cualquier choque local común a ambos grupos
educativos dentro de un departamento-trimestre, incluidas las llegadas de migrantes y
las olas de regularización del periodo de estudio, deja prácticamente sin cambios las
estimaciones de formalidad y de trabajo independiente ($-0.041$ y $+0.034$); solo la
estimación de empleo se atenúa aún más, en línea con mi lectura de ese margen. En
segundo lugar, las estimaciones son estables, de hecho algo mayores, cuando excluyo a
Lima y Callao, y cuando excluyo los ocho departamentos que albergan a la gran mayoría de
la población venezolana. En tercer lugar, la geografía va en la dirección equivocada
para que la migración sea un factor de confusión: los migrantes se concentran en
departamentos de baja incidencia, de modo que la presión de desplazamiento derivada de
la inmigración generaría efectos concentrados donde encuentro los más pequeños, lo
contrario de la relación dosis-respuesta observada en la exposición regional. El
calendario apunta en la misma dirección: el pico de entradas de 2018--2019 es anterior a
mi ventana, y la prueba prepandemia de la siguiente subsección, estimada precisamente en
esos años de máxima afluencia, muestra que el contraste educativo deriva en la dirección
opuesta a mis estimaciones de 2022.

\subsection{La crisis política de 2022--2023}

Un segundo choque agregado se superpone a la ventana posterior a la reforma: la crisis
política que siguió a la destitución del presidente Castillo el 7 de diciembre de 2022,
con protestas, bloqueos de carreteras y estados de emergencia entre diciembre de 2022 y
marzo de 2023, seguida de una recesión en 2023. Las protestas se concentraron en la
sierra sur, y varios de sus epicentros son también departamentos de alta incidencia, de
modo que la preocupación es específica: las caídas del empleo formal provocadas por el
conflicto podrían cargarse sobre el gradiente regional del índice de Kaitz, que sostiene
parte del peso causal de este artículo. Para el contraste educativo, los efectos fijos
departamento$\times$trimestre del Cuadro~\ref{tab:migration} ya absorben cualquier choque
departamento-trimestre, incluida la crisis; el diseño regional, en cambio, se identifica
precisamente a partir de la variación por departamento y en el tiempo, y requiere una
respuesta directa.

Mido la intensidad local de la crisis con los reportes mensuales de conflictos de la
Presidencia del Consejo de Ministros, que registran, para cada departamento, el número
máximo de personas movilizadas en acciones de protesta; tomo el pico de cada
departamento entre enero y febrero de 2023, los meses en que la movilización nacional
alcanzó su máximo, lo expreso por residente en edad de trabajar y lo estandarizo entre
departamentos (el Apéndice A detalla la fuente). Tres resultados del Cuadro del
Apéndice~\ref{tab:crisis} indican que la crisis no genera el gradiente regional. En
primer lugar, la geografía de las protestas no es la geografía de la incidencia: la
correlación entre departamentos de la intensidad de las protestas y la exposición según
el índice de Kaitz es de 0.19. Hu\'anuco, el departamento con el mayor índice de Kaitz,
casi no tuvo movilizaciones, mientras que Arequipa y Tacna, entre los más movilizados, se
ubican en la mitad inferior de la distribución de la exposición. En segundo lugar,
controlar por la intensidad de las protestas interactuada con el indicador del periodo
posterior, o con un indicador de 2023 en adelante, deja las pendientes del índice de
Kaitz prácticamente sin cambios (formalidad: $-0.018$ en la especificación base frente a
$-0.016$ a $-0.017$ con los controles, con $p$ por aleatorización $\le 0.001$ en todos
los casos; trabajo independiente: $+0.010$ frente a $+0.008$), mientras que el propio
control de protestas entra de manera significativa y con los signos esperados, de modo
que la medida tiene poder de identificación y sencillamente no explica el gradiente. En
tercer lugar, excluir los seis departamentos del sur que estuvieron en el centro de las
protestas deja intacto el gradiente ($-0.016$ para la formalidad, $p$ por aleatorización
$= 0.002$).

Un patrón temporal merece una declaración franca. Restringir la ventana posterior a la
reforma a 2022Q3--Q4, los dos trimestres posteriores a la reforma pero anteriores a la
crisis, reduce la pendiente regional de la formalidad a $-0.006$, que pierde
significancia; el gradiente, por tanto, se estima sobre todo a partir de 2023, cuando
coexiste con la crisis y la recesión. Tres observaciones me impiden leer el gradiente
como un artefacto de la crisis: el contraste educativo ya está plenamente presente en
esos dos trimestres previos a la crisis (formalidad $-0.029$, $p<0.01$, con la misma
magnitud bajo efectos fijos departamento$\times$trimestre que absorben todos los choques
locales de la crisis); los estudios de eventos muestran que el efecto de composición se
construye gradualmente a lo largo del primer año, de modo que una pendiente pequeña en
la ventana temprana es lo que implican las dinámicas y no una discrepancia; y el control
directo de protestas descrito arriba, que es exactamente la variable sobre la que se
cargaría un factor de confusión ligado a la crisis, no mueve nada. Aun así, la lectura
justa es que la corroboración regional es más sólida para la persistencia de mediano
plazo del efecto, y que el contraste educativo con absorción de choques locales es el
que aporta la evidencia sobre el calendario inmediato.

\subsection{Una extensión prepandemia de la prueba de tendencias paralelas}

Como mi ventana se abre en plena recuperación pospandemia, pruebo también el contraste
educativo en tiempos normales, con los módulos anuales de empleo de la ENAHO de
2018--2019 y la ventana tranquila 2018Q3--2019Q4 (posterior al aumento de la RMV de
abril de 2018 y anterior a cualquier otra reforma). El ejercicio es doblemente
informativo. Por un lado, aconseja humildad sobre el diseño educativo en niveles: las
pruebas conjuntas de las interacciones $\mathrm{Low}\times$trimestre rechazan la nula
para los salarios, el empleo y el trabajo independiente incluso en esta ventana
tranquila: en tiempos normales, los resultados de baja y alta educación oscilan uno
respecto del otro a bajas frecuencias, lo que es un argumento para apoyar el peso causal
en el diseño de exposición regional y no solo en el contraste educativo. Por otro lado,
la \emph{dirección} de la deriva en tiempos normales va en contra de mis hallazgos de
2022, no a su favor: una reforma placebo ubicada en 2019Q1 produce un ``efecto''
\emph{negativo} sobre el trabajo independiente ($-2.2$ puntos porcentuales) y uno
positivo y no significativo sobre la formalidad ($+0.8$ puntos), la imagen especular del
patrón posterior a la reforma. Cualquiera que sea la divergencia secular entre los
grupos, antes de la pandemia empujaba a los trabajadores de baja educación hacia la
formalidad y los alejaba del trabajo independiente, de modo que sesga hacia cero las
estimaciones de reasignación de 2022 en lugar de generarlas.

\subsection{Validación fuera de muestra: la reforma de 2025}

Por último, replico todo el diseño en torno a la segunda reforma, que elevó el salario
mínimo a 1130 soles el 1 de enero de 2025, usando los trimestres de 2024 y 2025. Esta
reforma es de una magnitud casi idéntica a la de 2022, pero ocurre en un entorno más
tranquilo y pospandemia. La columna (6) del Cuadro~\ref{tab:inference} reporta las
estimaciones. No aparece ninguna respuesta salarial ($0.002$, error estándar $0.011$):
en dos reformas ocurridas en condiciones macroeconómicas distintas, no hay evidencia de
que el salario mínimo haya elevado la remuneración relativa de los trabajadores poco
calificados. La reasignación, en cambio, no se replica: los efectos sobre la formalidad y
el empleo son nulos en torno a la reforma de 2025, y el efecto sobre el trabajo
independiente es pequeño y de signo opuesto. Tomo este resultado en serio en lugar de
buscarle una explicación que lo neutralice. Implica que la reasignación que documento en
torno a 2022 no es una respuesta estructural estable que siga a todo aumento del salario
mínimo; operó cuando el nuevo piso barrió una banda densa de remuneraciones formales,
una condición que el aumento de 2025 ya no cumplía. La Sección~\ref{sec:discussion}
desarrolla esta interpretación de dependencia del estado, examina directamente sus dos
variables de estado candidatas y discute las alternativas, incluida la posibilidad de
que dinámicas residuales de la recuperación contribuyan a las estimaciones de 2022.
